\section{Our Approach}
\label{sec:approach}

% Owner: Christopher Rodolph. Tracking: GitHub issue #26.
% Each paragraph below maps to the GitHub issues noted in its comment;
% the overall roadmap is issue #17 and project-docs/PLAN.md.

We implement Speculative Taint Tracking (STT)~\cite{Yu:STT:2019}, followed by
the two $\mu$STT mechanisms~\cite{Chen:uSTT:2025}, in the RTL of SonicBOOM
(BOOMv3)~\cite{Zhao:SonicBOOM:2020}, the out-of-order RV64GC core of the
Chipyard SoC framework~\cite{Amid:Chipyard:2020}. $\mu$STT has so far been
evaluated only in gem5 on ARM. Its released RTL is a standalone study of the
rename-stage circuit, not a working core. Our contribution is to build both
mechanisms into a complete, synthesizable RISC-V core and measure their cost.
The work splits into four tracks that each own part of the problem:
taint tracking at rename (A), delaying transmitters at issue (B), security and
performance evaluation (C), and hardware cost (D).

\subsection{STT in BOOM terms}
STT treats a load that executes speculatively as an \emph{access}
instruction. The register it writes is \emph{tainted}, and the taint spreads to
every instruction that reads a tainted register. A \emph{transmitter} is an
instruction whose timing can reveal its operands: in BOOM, loads and stores
(through their address) and branches and \texttt{jalr} (through their outcome).
A tainted transmitter may not execute until its taint is cleared. Taint is
cleared once the load that caused it passes the \emph{visibility point} (VP)
and can no longer be squashed. We use the Spectre threat model
(Section~\ref{sec:threat}), in which the VP is the oldest unresolved branch.
Each register's taint is stored as its \emph{youngest root of taint} (YRoT):
the ROB index of the youngest access load it depends on. An instruction is
safe once its YRoT is older than the VP.

\subsection{Building the mechanism}

% Issue #3 (A1): interface spec, shared by tracks A and B.
\paragraph{Interface}
Tracks A and B first agree on a shared interface. Each micro-op gains a YRoT
field and a transmitter flag. The ROB exports the VP, together with an
\textit{isSafe(YRoT)} check that handles ROB wraparound. New configuration
switches (\texttt{enableSTT}, \texttt{impedeUop}, \texttt{threatModel}) default
to off. With every switch off, the generated Verilog must be identical to
the unmodified baseline, so every measurement compares like with like.

% Issue #4 (A2): baseline STT taint propagation at rename.
\paragraph{Taint propagation (track A)}
At rename, a taint table with one entry per physical register, modeled on
BOOM's busy table, gives each instruction a YRoT: the youngest YRoT among its sources. A load's
destination takes the load's own ROB index. Instructions renamed in the same
cycle must see each other's taint, which requires a bypass between them. The
taint state must also be restored on branch mispredicts and exceptions, along
with the rename map table.

% Issue #6 (B1): conservative stand-in; #7 (B2): real STT delay + LSU audit.
\paragraph{Delaying transmitters (track B)}
Track B starts without waiting for track A. A \emph{conservative stand-in}
blocks every transmitter from issuing while any older branch is unresolved,
which BOOM already tracks in the branch mask of each micro-op. This policy is
strictly stronger than STT. It lets security testing begin early and gives an
upper bound on the performance cost. Track B then replaces the branch-mask
condition with track A's \textit{isSafe(YRoT)} check. Blocking issue at the issue
queue is not enough on its own, because BOOM's load/\allowbreak store unit can fire loads
without a fresh issue: load-queue retries, speculative wakeup, TLB-miss
replay, store-to-load forwarding, and the prefetcher. We audit each of these
paths so that all of them respect taint.

% Issue #5 (A3): Age Matrix; #8 (B3): impede micro-op.
\paragraph{The $\mu$STT mechanisms}
$\mu$STT addresses the two costs of this baseline. The \emph{Age Matrix}
replaces the cross-instruction age comparisons in the same-cycle taint bypass
with a shallow circuit, to cut rename delay. We port it to Chisel at BOOM's
widths and check that it produces exactly the same YRoTs as the baseline. The
\emph{impede micro-op} removes the per-entry delay logic from the issue
queue. A special micro-op sits in front of each transmitter and waits until
the VP passes its YRoT. The transmitter depends on it, so BOOM's ordinary
wakeup logic delays the transmitter. gem5's ARM model inserts this micro-op
by splitting instructions at decode, which BOOM does not do. We therefore
compare two options: emitting an extra micro-op at dispatch, or reserving a
dedicated issue-queue entry.

\subsection{Evaluation methodology}

% Issues #9 (C1), #10 (C2), #11 (C3).
\paragraph{Security and performance (track C)}
We port the BOOM Spectre proofs of concept (v1, v2, and return-stack-buffer
variants) to Chipyard bare-metal tests. Each must recover a secret on the
baseline and fail on every protected configuration. A regression script runs
the RISC-V ISA tests, the attacks, and optionally Spike co-simulation against
any configuration. Performance is measured as cycles and IPC from RTL
simulation in Verilator, using CoreMark, Dhrystone, and the riscv-tests
benchmarks, and normalized to the baseline.

% Issues #12 (D1), #13 (D2), #14 (D3), #15 (D4), #16 (D5).
\paragraph{Hardware cost (track D)}
We have no FPGA board, so hardware cost comes from Vivado 2023.2 targeting the
Xilinx VCU118 (\texttt{xcvu9p}), using Chipyard's existing BOOM FPGA
configuration. The full flow runs synthesis through implementation. It gives
hierarchical utilization (to attribute logic to rename, issue, ROB and LSU),
timing slack, and power estimates. Out-of-context synthesis of the BOOM tile
alone gives a faster loop for checking design changes. We also synthesize
$\mu$STT's released rename-stage RTL at BOOM's parameters, to reproduce the
paper's area and delay claims and give a reference for our Chisel versions.
Vivado simulation of the synthesized design runs short, targeted programs
such as the Spectre gadget.

\subsection{Putting it together}
Table~\ref{tab:configs} lists the configurations that combine these pieces.
Each is measured on all three axes (security, performance, and hardware cost)
against the same frozen baseline: unmodified Chipyard~1.15.0 with SonicBOOM.

\begin{table}[t]
  \caption{Configurations evaluated.}
  \label{tab:configs}
  \small
  \begin{tabular}{@{}lll@{}}
    \toprule
    Configuration & Taint (rename) & Delay (issue/LSU) \\
    \midrule
    Baseline      & none                & none \\
    Conservative  & none                & any unresolved branch \\
    STT           & taint table + bypass & \textit{isSafe(YRoT)} \\
    $\mu$STT      & Age Matrix          & impede micro-op \\
    \bottomrule
  \end{tabular}
\end{table}
